% Informe de laboratorio de Desarrollo de Videojuegos
% El punto de entrada organizado del informe es main.tex.

% Compilar con LuaLaTeX: fuentes TrueType reales de Windows.
\usepackage{fontspec}
\setmainfont{times.ttf}[Path=C:/Windows/Fonts/,BoldFont=timesbd.ttf,ItalicFont=timesi.ttf,BoldItalicFont=timesbi.ttf]
\setsansfont{times.ttf}[Path=C:/Windows/Fonts/,BoldFont=timesbd.ttf,ItalicFont=timesi.ttf,BoldItalicFont=timesbi.ttf]
\setmonofont{Consolas}[Scale=MatchLowercase]
% Compatibilidad de los enlaces internos con LuaTeX.
\protected\def\pdfinfo{\pdfextension info}
\protected\def\pdfdest{\pdfextension dest}
\protected\def\pdfstartlink{\pdfextension startlink}
\protected\def\pdfendlink{\pdfextension endlink}

\usepackage{graphicx}
\usepackage{xcolor}
\usepackage{array}
\usepackage{float}
\usepackage{tikz}
\usepackage{listings}
\usepackage{tcolorbox}
\tcbuselibrary{listings}
\usetikzlibrary{arrows.meta,positioning,shapes.geometric}

\definecolor{ESPEBlue}{HTML}{18345C}
\definecolor{ESPEGreen}{HTML}{08733F}
\definecolor{ESPEGreenDark}{HTML}{164D35}
\definecolor{PaleGreen}{HTML}{EAF3ED}
\definecolor{AppBlue}{HTML}{345DEE}
\definecolor{PaleBlue}{HTML}{EAF0FF}
\definecolor{SoftGray}{HTML}{F3F5F3}
\definecolor{Ink}{HTML}{243247}
\definecolor{OrangeAI}{HTML}{BC571F}

\setlength{\topmargin}{-0.3in}
\setlength{\oddsidemargin}{-0.1in}
\setlength{\textwidth}{6.3in}
\setlength{\textheight}{9.0in}
\setlength{\headheight}{15pt}
\setlength{\headsep}{0.2in}
\setlength{\footskip}{0.35in}
\setlength{\parskip}{0.45em}
\setlength{\parindent}{1.1em}
\setlength{\emergencystretch}{3em}
\setlength{\tabcolsep}{4.5pt}
\renewcommand{\arraystretch}{1.25}
\renewcommand{\topfraction}{0.92}
\renewcommand{\bottomfraction}{0.78}
\renewcommand{\textfraction}{0.06}
\renewcommand{\floatpagefraction}{0.62}
\setlength{\textfloatsep}{14pt plus 4pt minus 4pt}
\setlength{\floatsep}{14pt plus 4pt minus 4pt}
\pagestyle{headings}
\renewcommand{\chaptername}{Capítulo}
\renewcommand{\contentsname}{Índice general}
\renewcommand{\listfigurename}{Índice de figuras}
\renewcommand{\listtablename}{Índice de tablas}
\renewcommand{\figurename}{Figura}
\renewcommand{\tablename}{\textcolor{ESPEGreen}{Tabla}}
\renewcommand{\appendixname}{Anexo}
\renewcommand{\bibname}{Referencias}
% Tablas: Times New Roman, encabezado en banda verde y reglas sencillas.
\let\originaltabular\tabular
\let\endoriginaltabular\endtabular
\renewenvironment{tabular}[1]{\small\sffamily\color{Ink}\originaltabular{#1}}{\endoriginaltabular}
\newcommand{\thd}[1]{%
  \begingroup
  \setlength{\fboxsep}{3pt}%
  \colorbox{ESPEGreen}{%
    \parbox[t]{\dimexpr\linewidth-2\fboxsep\relax}{%
      \raggedright\sffamily\small\bfseries\color{white}%
      \rule{0pt}{2.6ex}#1\par}%
  }%
  \endgroup
}
\newcolumntype{P}[1]{>{\raggedright\arraybackslash}p{#1}}
\newcommand{\toprule}{\noalign{\vskip 2pt\color{ESPEGreen}\hrule height 1.1pt\vskip 2pt}}
\newcommand{\midrule}{\noalign{\vskip 1pt\color{ESPEGreen!40}\hrule height 0.45pt\vskip 1pt}}
\newcommand{\bottomrule}{\noalign{\vskip 2pt\color{ESPEGreen}\hrule height 1.1pt}}

\newcommand{\url}[1]{\texttt{\detokenize{#1}}}
\pdfinfo{/Title (Tic-Tac-Toe en .NET MAUI - Minimax) /Subject (Informe de laboratorio - Desarrollo de Videojuegos)}

% Enlaces internos del índice con primitivas pdfTeX, sin paquetes externos.
\newcounter{toclinkid}
\makeatletter
\renewcommand{\addcontentsline}[3]{%
  \stepcounter{toclinkid}%
  \edef\@currdest{toclink\arabic{toclinkid}}%
  \pdfdest name{\@currdest} xyz%
  \addtocontents{#1}{\protect\contentsline{#2}{#3}{\thepage}{\@currdest}}%
}
\renewcommand{\contentsline}[4]{%
  \csname l@#1\endcsname%
    {\pdfstartlink attr{/Border [0 0 0]} goto name{#4}%
     \textcolor{ESPEGreen}{#2}\pdfendlink}%
    {\pdfstartlink attr{/Border [0 0 0]} goto name{#4}%
     \textcolor{ESPEGreen}{#3}\pdfendlink}%
}
\makeatother

\setlength{\abovecaptionskip}{5pt}
\setlength{\belowcaptionskip}{3pt}
\makeatletter
\long\def\@makecaption#1#2{%
  \vskip\abovecaptionskip
  {\small\sffamily\centering
   \textcolor{ESPEGreen}{\bfseries #1}\hspace{0.5em}%
   \textcolor{Ink}{#2}\par}%
  \vskip\belowcaptionskip
}
\setlength{\@fptop}{0pt}
\setlength{\@fpbot}{0pt plus 1fil}
\makeatother

\makeatletter
\def\@makechapterhead#1{%
  \vspace*{22pt}%
  {\parindent\z@\raggedright\normalfont
    {\color{ESPEGreen}\sffamily\Large\bfseries\@chapapp\space\thechapter\par}%
    \vskip 8pt%
    {\color{Ink}\sffamily\bfseries\Huge #1\par}%
    \nobreak\vskip 22pt%
  }}
\def\@makeschapterhead#1{%
  \vspace*{22pt}%
  {\parindent\z@\raggedright\normalfont
    {\color{ESPEGreen}\sffamily\bfseries\Huge #1\par}%
    \nobreak\vskip 22pt%
  }}
\renewcommand{\chaptermark}[1]{\markboth{\sffamily\small\color{ESPEGreenDark}\chaptername\space\thechapter.\quad #1}{}}
\renewcommand{\sectionmark}[1]{\markright{\sffamily\small\color{ESPEGreenDark}\thesection\quad #1}}
\renewcommand{\section}{\@startsection{section}{1}{\z@}%
  {-3.4ex\@plus -1ex \@minus -.2ex}%
  {2ex \@plus .2ex}%
  {\normalfont\Large\sffamily\bfseries\color{ESPEGreenDark}}}
\renewcommand{\subsection}{\@startsection{subsection}{2}{\z@}%
  {-2.8ex\@plus -.8ex \@minus -.2ex}%
  {1.4ex \@plus .2ex}%
  {\normalfont\large\sffamily\bfseries\color{ESPEGreenDark}}}
\makeatother

% Dracula / cyberpunk: estilos de código centralizados.
\definecolor{DraculaBG}{HTML}{282A36}
\definecolor{DraculaPanel}{HTML}{21222C}
\definecolor{DraculaFG}{HTML}{F8F8F2}
\definecolor{DraculaPurple}{HTML}{BD93F9}
\definecolor{DraculaCyan}{HTML}{8BE9FD}
\definecolor{DraculaPink}{HTML}{FF79C6}
\definecolor{DraculaGreen}{HTML}{50FA7B}
\definecolor{DraculaComment}{HTML}{A0ACD7}

\lstdefinestyle{draculacyber}{
  language=[Sharp]C,
  basicstyle=\ttfamily\small\color{DraculaFG},
  keywordstyle=\bfseries\color{DraculaPink},
  stringstyle=\color{DraculaGreen},
  commentstyle=\itshape\color{DraculaComment},
  identifierstyle=\color{DraculaFG},
  morekeywords={Minimax,MAX,MIN,si,sino,para,retornar,funcion,profundidad},
  columns=fullflexible,
  keepspaces=true,
  breaklines=true,
  breakatwhitespace=false,
  showstringspaces=false,
  numbers=left,
  numberstyle=\ttfamily\scriptsize\color{DraculaComment},
  numbersep=9pt,
  stepnumber=1,
  frame=none,
  backgroundcolor=\color{DraculaBG},
  xleftmargin=1.8em,
  xrightmargin=0.3em,
  aboveskip=0.5em,
  belowskip=0.5em,
  tabsize=4
}
\lstset{style=draculacyber}

\tcbset{draculacyberbox/.style={
  colback=DraculaBG,
  colframe=DraculaPurple,
  colbacktitle=DraculaPanel,
  coltitle=DraculaCyan,
  boxrule=0.7pt,
  leftrule=2pt,
  arc=1.5mm,
  left=2mm,
  right=2mm,
  top=2mm,
  bottom=2mm,
  before skip=12pt,
  after skip=12pt,
  fonttitle=\bfseries,
  titlerule=0.5pt
}}
\newtcblisting{codigobox}[1]{
  draculacyberbox,
  listing only,
  title={#1},
  listing options={style=draculacyber}
}
\newcommand{\estado}[1]{\textsf{\small #1}}
\newcommand{\ruta}[1]{\texttt{#1}}
